% GENERATED FILE. Edit canonical lessons, then run npm run generate:books.
% canonical-source-hash: fnv1a64:71ca509fa605b734
% canonical-lessons: TA-C372-cattam, TA-C372-miru, TA-C372-kataippiti, TA-C372-tanti, TA-C372-inanku

\chapter{Laws and Rules}
\label{ch:ta-laws-and-rules}
\hlchaptermodality{372}

\begin{chapteropening}
\textbf{By the end of this chapter:} \emph{I can say that someone keeps or breaks a law.}

Complete the last lesson of chapter 372: I can say that someone keeps or breaks a law.
\end{chapteropening}

\section[\texorpdfstring{\ta{சட்டம்} (caṭṭam)}{சட்டம் (caṭṭam)}]{\ta{சட்டம்} (caṭṭam) — a law}

\label{lesson:TA-C372-cattam}



\begin{quote}
Before the new one: say the Tamil for a majority, then the Tamil for to oppress.
\end{quote}

\subsection*{You'll want to know: \ta{சட்டம்}}
\textbf{\ta{சட்டம்}} — \emph{caṭṭam} — \textquotedblleft{}a law\textquotedblright{}.

\begin{quote}
\ta{அரசு} \ta{ஒரு} \ta{புதிய} \ta{சட்டம்} \ta{கொண்டுவந்தது}.
\end{quote}

\begin{quote}
\emph{aracu oru putiya caṭṭam koṇṭuvantatu.}
\end{quote}

\begin{quote}
The government brought in a new law.
\end{quote}

Plural: \ta{சட்டங்கள்}, \emph{caṭṭaṅkaḷ}.

\ta{சட்டப்படி}, by law. A -\ta{அம்} noun makes its plural in -\ta{ங்கள்}.

\begin{cousinweb}
\textbf{\ta{சட்டம்}} (\emph{caṭṭam}) is native Dravidian: DEDR 2304 lists Tamil \emph{caṭṭam}, \textquotedblleft{}wooden frame; plan, model, rule, order, regulation\textquotedblright{}: a frame that holds things in place.
\end{cousinweb}

\begin{grammarlens}[title={What you've built}]
One more word for giving reasons and opinions.
\end{grammarlens}

\subsection*{Your turn}
\begin{itemize}
\raggedright
  \item \emph{Say it:} \emph{caṭṭam}
  \item \emph{Say it:} \emph{caṭṭam}, once more
  \item \emph{Say it:} \emph{aracu oru putiya caṭṭam koṇṭuvantatu.} — The government brought in a new law.
  \item \emph{Recall:} say the Tamil for a majority, then the Tamil for to oppress, then say \emph{caṭṭam} again
\end{itemize}

\begin{tcolorbox}[breakable,colback=teal!4,colframe=teal!35!black,title={Before you move on}]
What does \ta{சட்டம்} mean? (\textquotedblleft{}A law\textquotedblright{}.) Say it once more.
\end{tcolorbox}
\section[\texorpdfstring{\ta{மீறு} (mīṟu)}{மீறு (mīṟu)}]{\ta{மீறு} (mīṟu) — to violate; to break (a rule)}

\label{lesson:TA-C372-miru}



\begin{quote}
Before the new one: say the Tamil for to oppress, then the Tamil for a law.
\end{quote}

\subsection*{You'll want to know: \ta{மீறு}}
\textbf{\ta{மீறு}} — \emph{mīṟu} — \textquotedblleft{}to violate; to break (a rule)\textquotedblright{}.

\begin{quote}
\ta{அவன்} \ta{சாலை} \ta{விதிகளை} \ta{மீறினான்}.
\end{quote}

\begin{quote}
\emph{avaṉ cālai vitikaḷai mīṟiṉāṉ.}
\end{quote}

\begin{quote}
He broke the traffic rules.
\end{quote}

Past (I …): \ta{மீறினேன்}, \emph{mīṟiṉēṉ}.

What is broken takes -\ta{ஐ}: \ta{விதியை} \ta{மீறு}. The noun is \ta{மீறல்}, a violation.

\begin{cousinweb}
\textbf{\ta{மீறு}} (\emph{mīṟu}) is native Dravidian: DEDR 4884 lists Tamil \emph{mīṟu}, \textquotedblleft{}to go beyond, exceed, transgress, violate\textquotedblright{}.
\end{cousinweb}

\begin{grammarlens}[title={What you've built}]
One more word for giving reasons and opinions.
\end{grammarlens}

\subsection*{Your turn}
\begin{itemize}
\raggedright
  \item \emph{Say it:} \emph{mīṟu}
  \item \emph{Say it:} \emph{mīṟu}, once more
  \item \emph{Say it:} \emph{avaṉ cālai vitikaḷai mīṟiṉāṉ.} — He broke the traffic rules.
  \item \emph{Recall:} say the Tamil for to oppress, then the Tamil for a law, then say \emph{mīṟu} again
\end{itemize}

\begin{tcolorbox}[breakable,colback=teal!4,colframe=teal!35!black,title={Before you move on}]
What does \ta{மீறு} mean? (\textquotedblleft{}To violate; to break (a rule)\textquotedblright{}.) Say it once more.
\end{tcolorbox}
\section[\texorpdfstring{\ta{கடைப்பிடி} (kaṭaippiṭi)}{கடைப்பிடி (kaṭaippiṭi)}]{\ta{கடைப்பிடி} (kaṭaippiṭi) — to observe; to abide by}

\label{lesson:TA-C372-kataippiti}



\begin{quote}
Before the new one: say the Tamil for a law, then the Tamil for to violate.
\end{quote}

\subsection*{You'll want to know: \ta{கடைப்பிடி}}
\textbf{\ta{கடைப்பிடி}} — \emph{kaṭaippiṭi} — \textquotedblleft{}to observe; to abide by\textquotedblright{}.

\begin{quote}
\ta{நாங்கள்} \ta{பள்ளி} \ta{விதிகளைக்} \ta{கடைப்பிடிக்கிறோம்}.
\end{quote}

\begin{quote}
\emph{nāṅkaḷ paḷḷi vitikaḷaik kaṭaippiṭikkiṟōm.}
\end{quote}

\begin{quote}
We abide by the school rules.
\end{quote}

Past (I …): \ta{கடைப்பிடித்தேன்}, \emph{kaṭaippiṭittēṉ}.

What is kept takes -\ta{ஐ}; customs too: \ta{பழக்கத்தைக்} \ta{கடைப்பிடி}.

\begin{cousinweb}
\textbf{\ta{கடைப்பிடி}} (\emph{kaṭaippiṭi}) is built from native \emph{kaṭai}, \textquotedblleft{}end, limit\textquotedblright{} (DEDR 1109), and \emph{piṭi}, \textquotedblleft{}to catch, hold\textquotedblright{} (DEDR 4148): to hold on to the end.
\end{cousinweb}

\begin{grammarlens}[title={What you've built}]
One more word for giving reasons and opinions.
\end{grammarlens}

\subsection*{Your turn}
\begin{itemize}
\raggedright
  \item \emph{Say it:} \emph{kaṭaippiṭi}
  \item \emph{Say it:} \emph{kaṭaippiṭi}, once more
  \item \emph{Say it:} \emph{nāṅkaḷ paḷḷi vitikaḷaik kaṭaippiṭikkiṟōm.} — We abide by the school rules.
  \item \emph{Recall:} say the Tamil for a law, then the Tamil for to violate, then say \emph{kaṭaippiṭi} again
\end{itemize}

\begin{tcolorbox}[breakable,colback=teal!4,colframe=teal!35!black,title={Before you move on}]
What does \ta{கடைப்பிடி} mean? (\textquotedblleft{}To observe; to abide by\textquotedblright{}.) Say it once more.
\end{tcolorbox}
\section[\texorpdfstring{\ta{தண்டி} (taṇṭi)}{தண்டி (taṇṭi)}]{\ta{தண்டி} (taṇṭi) — to punish}

\label{lesson:TA-C372-tanti}



\begin{quote}
Before the new one: say the Tamil for to violate, then the Tamil for to observe.
\end{quote}

\subsection*{You'll want to know: \ta{தண்டி}}
\textbf{\ta{தண்டி}} — \emph{taṇṭi} — \textquotedblleft{}to punish\textquotedblright{}.

\begin{quote}
\ta{தவறு} \ta{செய்தவர்களை} \ta{அரசு} \ta{தண்டிக்க} \ta{வேண்டும்}.
\end{quote}

\begin{quote}
\emph{tavaṟu ceytavarkaḷai aracu taṇṭikka vēṇṭum.}
\end{quote}

\begin{quote}
The government must punish those who did wrong.
\end{quote}

Past (I …): \ta{தண்டித்தேன்}, \emph{taṇṭittēṉ}.

Who is punished takes -\ta{ஐ}. The noun is \ta{தண்டனை}, a punishment.

\begin{cousinweb}
\textbf{\ta{தண்டி}} (\emph{taṇṭi}) is a verb built on Sanskrit \emph{daṇḍa}, \textquotedblleft{}stick, club; control, punishment\textquotedblright{} (CDIAL 6128). Tamil writes the d as \ta{த}.
\end{cousinweb}

\begin{grammarlens}[title={What you've built}]
One more word for giving reasons and opinions.
\end{grammarlens}

\subsection*{Your turn}
\begin{itemize}
\raggedright
  \item \emph{Say it:} \emph{taṇṭi}
  \item \emph{Say it:} \emph{taṇṭi}, once more
  \item \emph{Say it:} \emph{tavaṟu ceytavarkaḷai aracu taṇṭikka vēṇṭum.} — The government must punish those who did wrong.
  \item \emph{Recall:} say the Tamil for to violate, then the Tamil for to observe, then say \emph{taṇṭi} again
\end{itemize}

\begin{tcolorbox}[breakable,colback=teal!4,colframe=teal!35!black,title={Before you move on}]
What does \ta{தண்டி} mean? (\textquotedblleft{}To punish\textquotedblright{}.) Say it once more.
\end{tcolorbox}
\section[\texorpdfstring{\ta{இணங்கு} (iṇaṅku)}{இணங்கு (iṇaṅku)}]{\ta{இணங்கு} (iṇaṅku) — to comply; to give in}

\label{lesson:TA-C372-inanku}



\begin{quote}
Before the new one: say the Tamil for to observe, then the Tamil for to punish.
\end{quote}

\subsection*{You'll want to know: \ta{இணங்கு}}
\textbf{\ta{இணங்கு}} — \emph{iṇaṅku} — \textquotedblleft{}to comply; to give in\textquotedblright{}.

\begin{quote}
\ta{கடைசியில்} \ta{அவன்} \ta{எங்கள்} \ta{கோரிக்கைக்கு} \ta{இணங்கினான்}.
\end{quote}

\begin{quote}
\emph{kaṭaiciyil avaṉ eṅkaḷ kōrikkaikku iṇaṅkiṉāṉ.}
\end{quote}

\begin{quote}
In the end he gave in to our request.
\end{quote}

Past (I …): \ta{இணங்கினேன்}, \emph{iṇaṅkiṉēṉ}.

With -\ta{க்கு}: \ta{கோரிக்கைக்கு} \ta{இணங்கு}, to agree to a request.

\begin{cousinweb}
\textbf{\ta{இணங்கு}} (\emph{iṇaṅku}) is native Dravidian: DEDR 457 lists Tamil \emph{iṇai}, \textquotedblleft{}to join, unite, agree\textquotedblright{}, and \emph{iṇaṅku}, \textquotedblleft{}to agree, be friendly\textquotedblright{}.
\end{cousinweb}

\begin{grammarlens}[title={What you've built}]
One more word for giving reasons and opinions.
\end{grammarlens}

\subsection*{Your turn}
\begin{itemize}
\raggedright
  \item \emph{Say it:} \emph{iṇaṅku}
  \item \emph{Say it:} \emph{iṇaṅku}, once more
  \item \emph{Say it:} \emph{kaṭaiciyil avaṉ eṅkaḷ kōrikkaikku iṇaṅkiṉāṉ.} — In the end he gave in to our request.
  \item \emph{Recall:} say the Tamil for to observe, then the Tamil for to punish, then say \emph{iṇaṅku} again
\end{itemize}

\begin{tcolorbox}[breakable,colback=teal!4,colframe=teal!35!black,title={Before you move on}]
What does \ta{இணங்கு} mean? (\textquotedblleft{}To comply; to give in\textquotedblright{}.) Say it once more.
\end{tcolorbox}
